With the rise of cloud computing and the \ac{IoT}, modern applications have been
revolutionized from monolithical units running on local premises to a multitude
of smaller, interconnected services deployed on shared infrastructures used by
multiple tenants. Sharing resources among multiple applications improves
scalability and efficiency, reducing costs at the same time. However, it also
poses novel security challenges compared to traditional systems.

Protecting the confidentiality, integrity and authenticity of data and
computations is paramount to ensure the correct operation of an application.
While there already exist established techniques to protect data at rest and in
transit, protecting data while in use is still an open research area. \acp{TEE}
tackle this problem by relying on hardware-based mechanisms to isolate code and
data during execution, leveraging extensions to the underlying computing
architecture to provide strong isolation and access control both in the main
processor and in memory. Code and data running inside a \ac{TEE} are then
shielded from the rest of the system, and even privileged software such as the
hypervisor is not allowed to access the secure environment, thereby
significantly reducing the overall attack surface.

In this Ph.D.\ dissertation, we analyze \acp{TEE} from a practical standpoint,
focusing on how to securely and effectively integrate them into novel or
existing applications and protocols, taking into account possible shortcomings
and limitations of this technology. First, we propose a framework to facilitate
the development and management of distributed applications in shared,
heterogeneous infrastructures. To do so, we employ process-based \acp{TEE} such
as Intel \ac{SGX} and Sancus to minimize the attack surface, leveraging secure
I/O capabilities to guarantee strong end-to-end integrity of application data.
To showcase the effectiveness of our work, we develop and evaluate a prototype
based on a smart home system.

As a second step, we study recent architectures such as AMD \sevsnp{} and Intel
\ac{TDX}, which allow running fully-fledged virtual machines inside a \ac{TEE}.
While this significantly reduces the development effort, it also causes an
increased attack surface and additional operational complexity, since a virtual
machine comprises several software components that are executed in sequence
during boot. Thus, it is essential to verify each stage of the boot process to
ensure that the virtual machine is set up correctly. However, our findings
reveal that, in current cloud settings, users have limited knowledge about the
software executing within the \ac{TEE}, since cloud providers have significant
control over the boot process of virtual machines and could potentially access
sensitive user data stored inside the \ac{TEE}. Therefore, it is essential to
clarify trust relationships between cloud providers and customers when using
\acp{TEE} in the cloud.

Third, we investigate the interactions between an application running within a
\ac{TEE} and the underlying untrusted environment. In particular, we examine the
problem of getting a stable and reliable notion of time inside a \ac{TEE},
assessing the capabilities of various architectures to provide a secure time
source and evaluating common use cases and applications that rely on time. Then,
we use this knowledge to propose a novel revocation scheme in the context of
\ac{V2X} credential management, which relies on \acp{TEE} embedded inside
vehicles. Our formal verification work shows that our scheme gives a guaranteed
upper bound on revocation time, and our evaluation highlights the efficiency and
scalability of our solution.

In summary, this dissertation shows how \acp{TEE} can be utilized to enhance the
security of existing applications and design novel solutions that can provide
additional benefits compared to traditional approaches. Furthermore, we also
raise the attention to non-obvious practical challenges that arise when
deploying \ac{TEE} workloads to untrusted infrastructures such as public clouds.